% GENERATED FILE. Edit canonical lessons, then run npm run generate:books.
% canonical-source-hash: fnv1a64:6a2cf72e9de48eb0
% canonical-lessons: RU-C252-khokkey, RU-C252-tennis, RU-C252-basketbol, RU-C252-beg, RU-C252-plavaniye

\chapter{Hockey, Tennis, Basketball, Running, Swimming}
\label{ch:ru-hockey-tennis-basketball-running-swimming}
\hlchaptermodality{252}

\begin{chapteropening}
\textbf{By the end of this chapter:} \emph{I can say hockey, tennis, basketball, running and swimming.}

Complete the last lesson of chapter 252: I can say hockey, tennis, basketball, running and swimming.
\end{chapteropening}

\section[\texorpdfstring{\ru{хоккей} (khokkéy)}{хоккей (khokkéy)}]{\ru{хоккей} (khokkéy) — hockey}

\label{lesson:RU-C252-khokkey}



\begin{quote}
Before the new one: say the Russian for a pigeon, then the Russian for a sparrow.
\end{quote}

\subsection*{You'll want to know: \ru{хоккей}}
\textbf{\ru{хоккей}} (\emph{khokkéy}) — \textquotedblleft{}hockey\textquotedblright{}.

\begin{cousinweb}
\textbf{\ru{хоккей}} comes from English \emph{hockey}.
\end{cousinweb}

\begin{grammarlens}[title={What you've built}]
One more word for this chapter.
\end{grammarlens}

\subsection*{Your turn}
\begin{itemize}
\raggedright
  \item \emph{Say it:} \emph{khokkéy}
  \item \emph{Say it:} \emph{khokkéy}, once more
  \item \emph{Say it:} say \emph{khokkéy}
  \item \emph{Recall:} say the Russian for a pigeon, then the Russian for a sparrow, then say \emph{khokkéy} again
\end{itemize}

\begin{tcolorbox}[breakable,colback=teal!4,colframe=teal!35!black,title={Before you move on}]
What does \ru{хоккей} mean? (\textquotedblleft{}Hockey\textquotedblright{}.) Say it once more.
\end{tcolorbox}
\section[\texorpdfstring{\ru{теннис} (ténnis)}{теннис (ténnis)}]{\ru{теннис} (ténnis) — tennis}

\label{lesson:RU-C252-tennis}



\begin{quote}
Before the new one: say the Russian for a sparrow, then the Russian for hockey.
\end{quote}

\subsection*{You'll want to know: \ru{теннис}}
\textbf{\ru{теннис}} (\emph{ténnis}) — \textquotedblleft{}tennis\textquotedblright{}.

\begin{cousinweb}
\textbf{\ru{теннис}} comes from English \emph{tennis} (Vasmer), which goes back to French \emph{tenez}, \textquotedblleft{}hold!\textquotedblright{}, called out by the server.
\end{cousinweb}

\begin{grammarlens}[title={What you've built}]
One more word for this chapter.
\end{grammarlens}

\subsection*{Your turn}
\begin{itemize}
\raggedright
  \item \emph{Say it:} \emph{ténnis}
  \item \emph{Say it:} \emph{ténnis}, once more
  \item \emph{Say it:} say \emph{ténnis}
  \item \emph{Recall:} say the Russian for a sparrow, then the Russian for hockey, then say \emph{ténnis} again
\end{itemize}

\begin{tcolorbox}[breakable,colback=teal!4,colframe=teal!35!black,title={Before you move on}]
What does \ru{теннис} mean? (\textquotedblleft{}Tennis\textquotedblright{}.) Say it once more.
\end{tcolorbox}
\section[\texorpdfstring{\ru{баскетбол} (basketból)}{баскетбол (basketból)}]{\ru{баскетбол} (basketból) — basketball}

\label{lesson:RU-C252-basketbol}



\begin{quote}
Before the new one: say the Russian for hockey, then the Russian for tennis.
\end{quote}

\subsection*{You'll want to know: \ru{баскетбол}}
\textbf{\ru{баскетбол}} (\emph{basketból}) — \textquotedblleft{}basketball\textquotedblright{}.

\begin{cousinweb}
\textbf{\ru{баскетбол}} comes from English \emph{basketball}, \textquotedblleft{}basket ball\textquotedblright{}.
\end{cousinweb}

\begin{grammarlens}[title={What you've built}]
One more word for this chapter.
\end{grammarlens}

\subsection*{Your turn}
\begin{itemize}
\raggedright
  \item \emph{Say it:} \emph{basketból}
  \item \emph{Say it:} \emph{basketból}, once more
  \item \emph{Say it:} say \emph{basketból}
  \item \emph{Recall:} say the Russian for hockey, then the Russian for tennis, then say \emph{basketból} again
\end{itemize}

\begin{tcolorbox}[breakable,colback=teal!4,colframe=teal!35!black,title={Before you move on}]
What does \ru{баскетбол} mean? (\textquotedblleft{}Basketball\textquotedblright{}.) Say it once more.
\end{tcolorbox}
\section[\texorpdfstring{\ru{бег} (beg)}{бег (beg)}]{\ru{бег} (beg) — running}

\label{lesson:RU-C252-beg}



\begin{quote}
Before the new one: say the Russian for tennis, then the Russian for basketball.
\end{quote}

\subsection*{You'll want to know: \ru{бег}}
\textbf{\ru{бег}} (\emph{beg}) — \textquotedblleft{}running\textquotedblright{}.

\begin{cousinweb}
\textbf{\ru{бег}} $\leftarrow$ Proto-Slavic \emph{*běgŭ}, with \emph{\ru{бежать}}, \textquotedblleft{}to run\textquotedblright{} — inherited. Vasmer relates it to Lithuanian \emph{bėgas}, \textquotedblleft{}a running, flight\textquotedblright{}.
\end{cousinweb}

\begin{grammarlens}[title={What you've built}]
One more word for this chapter.
\end{grammarlens}

\subsection*{Your turn}
\begin{itemize}
\raggedright
  \item \emph{Say it:} \emph{beg}
  \item \emph{Say it:} \emph{beg}, once more
  \item \emph{Say it:} say \emph{beg}
  \item \emph{Recall:} say the Russian for tennis, then the Russian for basketball, then say \emph{beg} again
\end{itemize}

\begin{tcolorbox}[breakable,colback=teal!4,colframe=teal!35!black,title={Before you move on}]
What does \ru{бег} mean? (\textquotedblleft{}Running\textquotedblright{}.) Say it once more.
\end{tcolorbox}
\section[\texorpdfstring{\ru{плавание} (plávaniye)}{плавание (plávaniye)}]{\ru{плавание} (plávaniye) — swimming}

\label{lesson:RU-C252-plavaniye}



\begin{quote}
Before the new one: say the Russian for basketball, then the Russian for running.
\end{quote}

\subsection*{You'll want to know: \ru{плавание}}
\textbf{\ru{плавание}} (\emph{plávaniye}) — \textquotedblleft{}swimming\textquotedblright{}.

\begin{cousinweb}
\textbf{\ru{плавание}} is a noun from \emph{\ru{плавать}}, \textquotedblleft{}to swim\textquotedblright{}, on the root of \emph{\ru{плыть}}, \textquotedblleft{}to float\textquotedblright{}. Vasmer relates \emph{\ru{плыть}} to Sanskrit \emph{plávatē}, \textquotedblleft{}floats\textquotedblright{}, and Greek \emph{pléō}, \textquotedblleft{}I sail\textquotedblright{}.
\end{cousinweb}

\begin{grammarlens}[title={What you've built}]
One more word for this chapter.
\end{grammarlens}

\subsection*{Your turn}
\begin{itemize}
\raggedright
  \item \emph{Say it:} \emph{plávaniye}
  \item \emph{Say it:} \emph{plávaniye}, once more
  \item \emph{Say it:} say \emph{plávaniye}
  \item \emph{Recall:} say the Russian for basketball, then the Russian for running, then say \emph{plávaniye} again
\end{itemize}

\begin{tcolorbox}[breakable,colback=teal!4,colframe=teal!35!black,title={Before you move on}]
What does \ru{плавание} mean? (\textquotedblleft{}Swimming\textquotedblright{}.) Say it once more.
\end{tcolorbox}
